% GENERATED FILE. Edit canonical lessons, then run npm run generate:books.
% canonical-source-hash: fnv1a64:32d1f62b30e1600e
% canonical-lessons: ZH-C99-xi, ZH-W99-xi-wash, ZH-C99-shuaya, ZH-C99-xilian, ZH-C99-xizao, ZH-C99-shutou

\chapter{Getting Washed}
\label{ch:zh-getting-washed}
\hlchaptermodality{99}

\begin{chapteropening}
\textbf{By the end of this chapter:} \emph{I can say wash, brush my teeth, wash my face, have a shower and comb my hair.}

Complete the last lesson of chapter 99: I can say wash, brush my teeth, wash my face, have a shower and comb my hair.
\end{chapteropening}

\section[\texorpdfstring{\zh{洗} (xǐ)}{洗 (xǐ)}]{\zh{洗} (xǐ) — to wash}

\label{lesson:ZH-C99-xi}



\begin{quote}
Before the new one: say the Mandarin for every day, then the Mandarin for after.
\end{quote}

\subsection*{You'll want to know: \zh{洗}}
\textbf{\zh{洗}} — \emph{xǐ} — \textquotedblleft{}to wash\textquotedblright{}.

\begin{quote}
\textbf{\zh{我去洗手}} \emph{wǒ qù xǐ shǒu} — I'm going to wash my hands.
\end{quote}

The new character, in pieces: \zh{洗} — the water piece of \zh{汉} on the left; on the right a six-stroke piece read xiān, the first syllable of xiānsheng. The next lesson puts it in your hand.

\begin{cousinweb}
\textbf{\zh{洗}} (\emph{xǐ}) is inherited from Middle Chinese \emph{sejX} (Guangyun, in Baxter's transcription), glossed 'to wash and bathe'.
\end{cousinweb}

\begin{grammarlens}[title={What you've built}]
One more word for everyday actions.
\end{grammarlens}

\subsection*{Your turn}
\begin{itemize}
\raggedright
  \item \emph{Say it:} \emph{xǐ}
  \item \emph{Say it:} \emph{xǐ}, once more
  \item \emph{Say it:} \emph{wǒ qù xǐ shǒu} — I'm going to wash my hands.
  \item \emph{Recall:} say the Mandarin for every day, then the Mandarin for after, then say \emph{xǐ} again
\end{itemize}

\begin{tcolorbox}[breakable,colback=teal!4,colframe=teal!35!black,title={Before you move on}]
What does \zh{洗} mean? (\textquotedblleft{}To wash\textquotedblright{}.) Say it once more.
\end{tcolorbox}
\section[\texorpdfstring{\zh{洗} (xǐ)}{洗 (xǐ)}]{\zh{洗} — the water piece, then xiān}

\label{lesson:ZH-W99-xi-wash}



\begin{quote}
The word first, then the character inside it.

\begin{itemize}
\raggedright
  \item \emph{Recall:} say \emph{xǐ} — \textbf{R1}, one lesson back, before it decays
\end{itemize}

\begin{quote}
\zh{汉} \zh{没} \zh{儿}
\end{quote}

\emph{Read it:} each one aloud, with nothing said first — \emph{hàn}, \emph{méi}, \emph{ér}
\end{quote}

\subsection*{Script — the water piece, then xiān}
\hlblockfigure{\detokenize{figures/ZH-W99-xi-wash-filmstrip.pdf}}{How \zh{洗} is written, stroke by stroke}

\begin{quote}
\zh{洗}
\end{quote}

\textbf{Two pieces, left before right.} Count the pieces, not the strokes.

\textbf{The water piece, three strokes}, as in \zh{汉} and \zh{没}: a dot, a second dot under it, then a stroke that rises from the bottom left.

\textbf{The right piece, xiān, six strokes}: a short left-falling stroke, a horizontal, a vertical down through it, a long horizontal; then a left-falling stroke, and a vertical that bends right along the bottom and hooks up, as at the foot of \zh{儿}.

3 + 6 = nine.

\subsection*{Your turn}
\begin{itemize}
\raggedright
  \item \emph{Write it:} \zh{洗}, all nine, with the model in front of you
  \item \emph{Cover:} write it again from memory, and count the pieces, not the strokes
\end{itemize}

\begin{tcolorbox}[breakable,colback=teal!4,colframe=teal!35!black,title={Before you move on}]
How many pieces in \zh{洗}? (\textbf{Two} — the water piece, then xiān.) Which characters you know have the same water piece? (\textbf{\zh{汉}} and \textbf{\zh{没}}.) What does the last stroke look like? (\textbf{The right stroke of \zh{儿}}.)

Source: \href{https://raw.githubusercontent.com/chanind/hanzi-writer-data/68d10a4b21150cae5e1ebbd223eed289cf32d90c/data/\%E6\%B4\%97.json}{Hanzi Writer Data \zh{洗}, pinned stroke paths and medians}.
\end{tcolorbox}
\section[shuāyá]{shuāyá — to brush one's teeth}

\label{lesson:ZH-C99-shuaya}



\begin{quote}
Before the new one: say the Mandarin for after, then the Mandarin for to wash.
\end{quote}

\subsection*{You'll want to know: shuāyá}
\emph{shuāyá} — \textquotedblleft{}to brush one's teeth\textquotedblright{}.

\begin{quote}
\emph{wǒ měitiān zǎoshang shuāyá} — I brush my teeth every morning.
\end{quote}

Its characters come later; for now it is a word for the ear.

\begin{cousinweb}
\emph{shuāyá} is built from shuā 'to brush' and yá 'tooth', both inherited from Middle Chinese, \emph{srwaet} and \emph{ngae} (Guangyun, in Baxter's transcription), glossed 'to brush and wipe' and 'teeth'.
\end{cousinweb}

\begin{grammarlens}[title={What you've built}]
One more word for everyday actions.
\end{grammarlens}

\subsection*{Your turn}
\begin{itemize}
\raggedright
  \item \emph{Say it:} \emph{shuāyá}
  \item \emph{Say it:} \emph{shuāyá}, once more
  \item \emph{Say it:} \emph{wǒ měitiān zǎoshang shuāyá} — I brush my teeth every morning.
  \item \emph{Recall:} say the Mandarin for after, then the Mandarin for to wash, then say \emph{shuāyá} again
\end{itemize}

\begin{tcolorbox}[breakable,colback=teal!4,colframe=teal!35!black,title={Before you move on}]
What does shuāyá mean? (\textquotedblleft{}To brush one's teeth\textquotedblright{}.) Say it once more.
\end{tcolorbox}
\section[xǐliǎn]{xǐliǎn — to wash one's face}

\label{lesson:ZH-C99-xilian}



\begin{quote}
Before the new one: say the Mandarin for to wash, then the Mandarin for to brush one's teeth.
\end{quote}

\subsection*{You'll want to know: xǐliǎn}
\emph{xǐliǎn} — \textquotedblleft{}to wash one's face\textquotedblright{}.

\begin{quote}
\emph{nǐ xǐliǎn le ma} — Have you washed your face?
\end{quote}

Its characters come later; for now it is a word for the ear.

Two third tones in a row: xǐ is said with a rising tone before liǎn, though both are written with the third tone.

\begin{cousinweb}
\emph{xǐliǎn} is built from \zh{洗} 'to wash' and liǎn 'face'. The first is inherited from Middle Chinese \emph{sejX} (Guangyun, in Baxter's transcription); the dictionary glosses the second, \emph{leamX}, as the name of a kind of stew, so how it came to mean face is not settled.
\end{cousinweb}

\begin{grammarlens}[title={What you've built}]
One more word for everyday actions.
\end{grammarlens}

\subsection*{Your turn}
\begin{itemize}
\raggedright
  \item \emph{Say it:} \emph{xǐliǎn}
  \item \emph{Say it:} \emph{xǐliǎn}, once more
  \item \emph{Say it:} \emph{nǐ xǐliǎn le ma} — Have you washed your face?
  \item \emph{Recall:} say the Mandarin for to wash, then the Mandarin for to brush one's teeth, then say \emph{xǐliǎn} again
\end{itemize}

\begin{tcolorbox}[breakable,colback=teal!4,colframe=teal!35!black,title={Before you move on}]
What does xǐliǎn mean? (\textquotedblleft{}To wash one's face\textquotedblright{}.) Say it once more.
\end{tcolorbox}
\section[xǐzǎo]{xǐzǎo — to have a shower; to have a bath}

\label{lesson:ZH-C99-xizao}



\begin{quote}
Before the new one: say the Mandarin for to brush one's teeth, then the Mandarin for to wash one's face.
\end{quote}

\subsection*{You'll want to know: xǐzǎo}
\emph{xǐzǎo} — \textquotedblleft{}to have a shower; to have a bath\textquotedblright{}.

\begin{quote}
\emph{wǒ wǎnshang xǐzǎo} — I have a shower in the evening.
\end{quote}

Its characters come later; for now it is a word for the ear.

Two third tones again: xǐ rises before zǎo, and is still written xǐ.

\begin{cousinweb}
\emph{xǐzǎo} is built from \zh{洗} and zǎo, both inherited from Middle Chinese, \emph{sejX} and \emph{tsawX} (Guangyun, in Baxter's transcription); the dictionary glosses the first 'to wash and bathe' and the second with the first: to wash.
\end{cousinweb}

\begin{grammarlens}[title={What you've built}]
One more word for everyday actions.
\end{grammarlens}

\subsection*{Your turn}
\begin{itemize}
\raggedright
  \item \emph{Say it:} \emph{xǐzǎo}
  \item \emph{Say it:} \emph{xǐzǎo}, once more
  \item \emph{Say it:} \emph{wǒ wǎnshang xǐzǎo} — I have a shower in the evening.
  \item \emph{Recall:} say the Mandarin for to brush one's teeth, then the Mandarin for to wash one's face, then say \emph{xǐzǎo} again
\end{itemize}

\begin{tcolorbox}[breakable,colback=teal!4,colframe=teal!35!black,title={Before you move on}]
What does xǐzǎo mean? (\textquotedblleft{}To have a shower; to have a bath\textquotedblright{}.) Say it once more.
\end{tcolorbox}
\section[shūtóu]{shūtóu — to comb one's hair}

\label{lesson:ZH-C99-shutou}



\begin{quote}
Before the new one: say the Mandarin for to wash one's face, then the Mandarin for to have a shower.
\end{quote}

\subsection*{You'll want to know: shūtóu}
\emph{shūtóu} — \textquotedblleft{}to comb one's hair\textquotedblright{}.

\begin{quote}
\emph{wǒ zǎoshang shūtóu} — I comb my hair in the morning.
\end{quote}

Its characters come later; for now it is a word for the ear.

\begin{cousinweb}
\emph{shūtóu} is built from shū 'to comb' and tóu 'head', both inherited from Middle Chinese, \emph{srjo} and \emph{duw} (Guangyun, in Baxter's transcription); the dictionary glosses the first as a comb, for putting the hair in order.
\end{cousinweb}

\begin{grammarlens}[title={What you've built}]
One more word for everyday actions.
\end{grammarlens}

\subsection*{Your turn}
\begin{itemize}
\raggedright
  \item \emph{Say it:} \emph{shūtóu}
  \item \emph{Say it:} \emph{shūtóu}, once more
  \item \emph{Say it:} \emph{wǒ zǎoshang shūtóu} — I comb my hair in the morning.
  \item \emph{Recall:} say the Mandarin for to wash one's face, then the Mandarin for to have a shower, then say \emph{shūtóu} again
\end{itemize}

\begin{tcolorbox}[breakable,colback=teal!4,colframe=teal!35!black,title={Before you move on}]
What does shūtóu mean? (\textquotedblleft{}To comb one's hair\textquotedblright{}.) Say it once more.
\end{tcolorbox}
